\section{Cross product and normal vectors}

The dot product produces a number and measures directional agreement. In three-dimensional space there is another important operation: the cross product. It takes two vectors and produces a new vector perpendicular to both of them.

This operation is specific to three-dimensional vector geometry in the form used here. It will provide the missing link between vectors and equations of planes.

\subsection*{Definition in coordinates}

Let
\[
u=(u_1,u_2,u_3),
\qquad
v=(v_1,v_2,v_3).
\]
Their cross product is defined by
\[
u\times v
=
\bigl(
u_2v_3-u_3v_2,\,
u_3v_1-u_1v_3,\,
u_1v_2-u_2v_1
\bigr).
\]

\begin{definitionbox}
For vectors \(u,v\in\mathbb{R}^3\), the cross product \(u\times v\) is the vector
\[
u\times v
=
\begin{pmatrix}
u_2v_3-u_3v_2\\
u_3v_1-u_1v_3\\
u_1v_2-u_2v_1
\end{pmatrix}.
\]
It is perpendicular to both \(u\) and \(v\).
\end{definitionbox}

A common mnemonic writes this computation as
\[
u\times v=
\begin{vmatrix}
\mathbf{i}&\mathbf{j}&\mathbf{k}\\
u_1&u_2&u_3\\
v_1&v_2&v_3
\end{vmatrix},
\]
with the usual alternating signs in the expansion. The coordinate formula above is the definition; the determinant notation is only a compact memory aid.

\begin{examplebox}
Let
\[
u=(1,2,3),
\qquad
v=(2,-1,1).
\]
Then
\[
u\times v
=
\bigl(
2\cdot1-3(-1),\,
3\cdot2-1\cdot1,\,
1(-1)-2\cdot2
\bigr),
\]
so
\[
u\times v=(5,5,-5).
\]
We can check perpendicularity:
\[
(5,5,-5)\cdot(1,2,3)=5+10-15=0,
\]
and
\[
(5,5,-5)\cdot(2,-1,1)=10-5-5=0.
\]
Thus the resulting vector is perpendicular to both original vectors.
\end{examplebox}

\subsection*{Order and orientation}

The order of the vectors matters:
\[
v\times u=-(u\times v).
\]
Reversing the order reverses the direction of the resulting normal vector, but it does not change the plane perpendicular to that normal direction.

\begin{remarkbox}
The vectors \(u\times v\) and \(v\times u\) have opposite orientations. In plane equations, either one may be used as a normal vector because multiplying a normal vector by a non-zero scalar does not change the plane.
\end{remarkbox}

\subsection*{Parallel vectors and the zero cross product}

If \(u\) and \(v\) are parallel, they do not determine two independent directions. In that case,
\[
u\times v=0.
\]
Conversely, for non-zero vectors in \(\mathbb{R}^3\),
\[
u\times v=0
\]
implies that the vectors are parallel.

\begin{theorembox}
For \(u,v\in\mathbb{R}^3\),
\[
u\times v=0
\]
if and only if \(u\) and \(v\) are linearly dependent.
\end{theorembox}

This criterion will be useful when deciding whether three points are collinear and whether two direction vectors really span a plane.

\subsection*{Geometric size}

If \(\theta\) is the angle between non-zero vectors \(u\) and \(v\), then
\[
\lVert u\times v\rVert
=
\lVert u\rVert\lVert v\rVert\sin\theta.
\]
Therefore \(\lVert u\times v\rVert\) is the area of the parallelogram spanned by \(u\) and \(v\).

This formula also explains why the cross product is zero for parallel vectors: their angle is \(0\) or \(\pi\), so \(\sin\theta=0\).

\subsection*{Finding a normal vector to a plane}

Suppose a plane contains two non-parallel direction vectors \(u\) and \(v\). A normal vector must be perpendicular to both directions. Therefore
\[
n=u\times v
\]
is a natural choice.

\begin{examplebox}
Consider the three points
\[
A=(1,0,0),
\qquad
B=(0,1,0),
\qquad
C=(0,0,1).
\]
Two direction vectors in their plane are
\[
u=B-A=(-1,1,0),
\]
and
\[
v=C-A=(-1,0,1).
\]
Their cross product is
\[
u\times v
=
\bigl(
1\cdot1-0\cdot0,\,
0(-1)-(-1)\cdot1,\,
(-1)\cdot0-1(-1)
\bigr)
=
(1,1,1).
\]
Thus a normal vector is
\[
n=(1,1,1).
\]
Using the point \(A\), the plane equation is
\[
n\cdot(P-A)=0,
\]
which becomes
\[
(x-1)+y+z=0,
\]
or
\[
x+y+z=1.
\]
\end{examplebox}

The chain of ideas is now complete:
\[
\text{three points}
\longrightarrow
\text{two direction vectors}
\longrightarrow
\text{cross product}
\longrightarrow
\text{normal vector}
\longrightarrow
\text{plane equation}.
\]

\begin{summarybox}
When using the cross product, ask:
\begin{itemize}
\item Are the vectors in \(\mathbb{R}^3\)?
\item Do I need a vector perpendicular to two given directions?
\item Have I preserved the alternating signs in the coordinate formula?
\item Does reversing the order only reverse orientation?
\item Is the cross product zero, indicating dependent directions?
\item Can the result be used as a normal vector to a plane?
\end{itemize}
\end{summarybox}

The cross product turns two independent directions inside a plane into one direction perpendicular to the plane.
